\section{Introduction}

A logical gate acts on encoded quantum information.  A transversal
implementation acts independently on the physical subsystems, so an error on
one subsystem is not spread directly to another subsystem by the gate.  This
makes transversal gates a basic fault-tolerant primitive, but also a strongly
constrained one: no quantum code admits a universal set of transversal gates
\cite{EastinKnill2009}.  Determining the exact transversal logical group of a
code family is therefore a natural structural problem.

This paper solves that problem for a family of one-logical-qudit quantum MDS
codes.  Let \(C\leq\F_q^{2m}\) be a linear
\([2m,m,m+1]_q\) maximum-distance-separable (MDS) code.  Its equal-phase CSS
state is
\begin{equation}\label{eq:equal-phase-css-state}
 |\Psi_C\rangle=q^{-m/2}\sum_{c\in C}|c\rangle.
\end{equation}
The state is absolutely maximally entangled (AME).  If one tensor factor is
regarded as an input, the remaining factors define a
\([[2m-1,1,m]]_q\) quantum MDS encoder.

Throughout, a \emph{transversal} physical unitary has the fixed-coordinate,
site-dependent form
\[
 U_1\otimes\cdots\otimes U_{2m-1}.
\]
The factors may differ from one coordinate to another.  This convention is
strictly broader than the uniform form \(U^{\otimes n}\).  The distinction is
essential here: on one branch of the classification, a logical symplectic
matrix is realized at different sites through conjugations determined by
coordinatewise multiplier ratios.  Logical groups are projective, so global
scalar phases are divided out.  Permutations of physical coordinates are not
included unless stated explicitly.

The classification is controlled by a standard coding-theoretic object.  A
vector \(s\in\F_q^{2m}\) is a diagonal code-to-dual multiplier when
\(\diag(s)C\subseteq C^\perp\).  Write
\[
 \mathcal D(C,C^\perp)
 =\{s\in\F_q^{2m}:\diag(s)C\subseteq C^\perp\}.
\]
The MDS property makes this apparently large endomorphism problem rigid: the
space has dimension zero or one, and any nonzero multiplier rescales every
coordinate and maps \(C\) onto \(C^\perp\).  We call the latter case
\emph{diagonal isoduality}.

\begin{theorem}[Exact transversal logical group]
\label{thm:diagonal-isodual-transversal-group}
Let \(q\) be an odd prime, let \(m\geq2\), and let
\(C\leq\F_q^{2m}\) be a linear \([2m,m,m+1]_q\) MDS code.  Then
\(\dim\mathcal D(C,C^\perp)\) is zero or one, and the projective group of
logical unitaries with fixed-coordinate transversal implementations is
\[
\begin{cases}
 \F_q^2\rtimes T,
   &\dim\mathcal D(C,C^\perp)=0,\\
 \F_q^2\rtimes\mathrm{SL}_2(q),
   &\dim\mathcal D(C,C^\perp)=1,
\end{cases}
\qquad
T=\{\diag(a,a^{-1}):a\in\F_q^\times\}.
\]
The second case is equivalent to \(SC=C^\perp\) for a nonsingular diagonal
matrix \(S\).
\end{theorem}

Thus the exact transversal group is determined by one nullity computation on
the classical code.  The translation factor \(\F_q^2\) consists of logical
Paulis.  The linear factor is either the split torus \(T\) or all of
\(\mathrm{SL}_2(q)\); no intermediate subgroup occurs.

The proof has three steps.  First, represent each one-qudit Clifford by a
\(2\)-by-\(2\) symplectic matrix.  Preservation of the CSS stabilizer gives
four block equations involving \(C\) and \(C^\perp\).  Any nonzero
off-diagonal block is a diagonal multiplier from one code to the other.
Second, the MDS distance forces every such multiplier to have full support and
forces the multiplier space to be at most one-dimensional.  If the space is
zero, all local symplectic matrices are diagonal and the linear logical image
is \(T\).  If it is one-dimensional, its generator constructs upper and lower
unipotent matrices, which generate \(\mathrm{SL}_2(q)\).  Finally, the
imported rigidity theorem below shows that an arbitrary product-unitary
implementation is Clifford at every site, proving that the constructed groups
are exact.

The result fits into a broad study of transversal and automorphism-induced
logical gates.  Rains organizes coordinatewise symplectic symmetries through
code-invariance algebras \cite[Theorems~4 and~14]{RainsNonbinary1999}.
Compatible symplectic bases describe transversal Clifford structure for
self-dual CSS codes \cite[Theorem~1]{TansuwannontTakadaFujii2025}, while
uniform diagonal multiblock actions and endomorphism algebras constrain qubit
stabilizer codes \cite[Section~II and Theorem~6.1]{DasuBurton2025}.
Code automorphisms and ZX dualities give further Clifford constructions,
including the Pauli corrections required to obtain the intended logical
action \cite[Sections~III--IV]{SayginelEtAl2025}.  Under the more restrictive
uniform convention \(V_F^{\otimes n}\), Prakash and Singhal characterize
odd-prime one-logical-qudit codes with a complete transversal Clifford set
\cite[Section~IV.C]{PrakashSinghal2026}; Albert gives related normal forms for
binary CSS codes \cite[Sections~III--V]{Albert2026}.

What is missing from these results is the exact fixed-coordinate,
site-dependent group for the present MDS--CSS family.  Theorem
\ref{thm:diagonal-isodual-transversal-group} supplies it by reducing the
problem to \(\mathcal D(C,C^\perp)\).  The key MDS statement and both group
constructions are proved here.  We import only the following consequences of
the companion paper: product-unitary rigidity, the encoder version of that
rigidity, a stabilizer-character correction, and minimum-support transition
data.  Tan's general AME--QMDS correspondence relates state symmetries and
transversal encoder gates \cite[Theorems~3.5--3.6]{Tan2026}; the present
multiplier theorem computes the exact fixed-coordinate subgroup for this
family.

\begin{theorem}[Imported stabilizer-AME rigidity
\cite{RuddAMEStabilizerRigidity2026}]
\label{thm:lu-lc-rigidity}
For every prime power \(q\) and \(m\geq2\), every product-unitary
intertwiner between two stabilizer \(\AME(2m,q)\) states is Clifford on each
tensor factor.
\end{theorem}

\begin{corollary}[Imported transversal Clifford no-go
\cite{RuddAMEStabilizerRigidity2026}]
\label{cor:transversal-clifford}
Every tensor-product conversion between the associated stabilizer
\([[2m-1,1,m]]_q\) encoders is Clifford on every physical factor and on the
logical factor.  In particular, no code in this family admits a transversal
non-Clifford logical unitary.
\end{corollary}

We also use one state-level consequence.  A Clifford map of two stabilizer
label groups need not preserve the chosen stabilizer characters, but a Pauli
correction removes that discrepancy.
\begin{lemma}[Imported stabilizer-character correction
\cite{RuddAMEStabilizerRigidity2026}]
\label{lem:pauli-phase-correction}
For stabilizer states in the standard finite-field Weyl convention, every
product-Clifford map of their stabilizer label groups can be corrected by a
product Pauli to map one state ray to the other.
\end{lemma}

The six-coordinate applications also use the companion paper's
minimum-support transition theorem.  If \(L\) is the Pauli-label group of a
stabilizer \(\AME(2m,q)\) state, then the labels supported on any
\((m+1)\)-coordinate set project bijectively onto the Pauli-label plane at
each retained coordinate, and these supported subgroups span \(L\).
Consequently, transition maps between overlapping supports determine local
Clifford equivalence up to local trace-symplectic changes of basis.  We use
only this statement; its proof and the companion paper's quantitative results
are not repeated here.

The rest of the paper has two roles.  First, it interprets the all-length
criterion at six coordinates.  A linear \([6,3,4]_q\) MDS code is represented
by six points in \(\PP^2(\F_q)\), no three collinear.  In this model, diagonal
isoduality means that the dual code is obtained by independently rescaling the
six coordinates; projective geometers call the same condition fixed-label
Gale self-association.  For six points in the projective plane, it is
equivalent to lying on a conic.  An explicit non-GRS one-parameter family then
has a degree-eight invariant \(z\) that classifies projective and monomial-code
equivalence over odd fields and local-Clifford and local-unitary equivalence
over odd prime fields.  The Clebsch \([6,3,4]_{11}\) code is a worked
syndrome-geometric example.

Second, the paper records refinements and limitations of the classification.
The linear \(\mathrm{SL}_2(q)\) action has coherent Weil lifts, whereas the
affine translation subgroup has the usual Heisenberg central extension.
Coordinate permutations form a separate group-extension problem.  Fixed-copy
scalar contractions are generically constant on regular code families, while
operator-valued reduced states retain the local Weyl directions used by the
rigidity theorem.  The finite six-point calculations have exact certificate
replays, but no finite enumeration enters the all-length theorem.

Section~\ref{sec:dictionary} fixes the stabilizer and MDS--CSS conventions.
Section~\ref{sec:exact-logical} proves the multiplier theorem and exact group,
then treats the lift refinements.  Section~\ref{sec:logical} translates the
criterion into six-point geometry.  Section~\ref{sec:pencil} classifies the
one-parameter family, and Section~\ref{sec:invariants} gives the limitation on
fixed-copy scalar invariants.  The appendices contain the explicit scalar
separators, transport calculation, and coordinate-permutation extensions.
Section~\ref{sec:verification} states the computational and formal trust
boundaries.
